% ===== PRÉAMBULE CANONIQUE — à coller VERBATIM en tête de CHAQUE .tex =====
\documentclass[11pt,a4paper]{article}
\usepackage[utf8]{inputenc}\usepackage[T1]{fontenc}\usepackage{lmodern}
\usepackage[french]{babel}
\usepackage{amsmath,amssymb,mathtools}\usepackage{icomma}
\usepackage{siunitx}\sisetup{locale=FR}  % \qty{}{}, \si{}, \ang{}
\usepackage{xcolor}\usepackage{colortbl}
\usepackage{tikz}\usepackage{pgfplots}\pgfplotsset{compat=1.18}
\usepackage[siunitx,europeanresistors]{circuitikz} % circuits (élec.)
\usepackage[most]{tcolorbox}\usepackage{enumitem}\usepackage{array,booktabs}
\usepackage[a4paper,margin=2.2cm]{geometry}
\usetikzlibrary{arrows.meta,calc,angles,quotes,positioning,patterns,%
  backgrounds,shapes.geometric,decorations.pathmorphing,decorations.markings,intersections}
\definecolor{primary}{RGB}{222,49,99}\definecolor{primarydark}{RGB}{150,22,55}
\definecolor{defcol}{RGB}{0,114,178}\definecolor{methcol}{RGB}{52,92,148}
\definecolor{excol}{RGB}{0,135,90}\definecolor{attncol}{RGB}{200,40,55}
\tcbset{boxbase/.style={enhanced,breakable,boxrule=0.9pt,arc=2.5pt,
  left=3mm,right=3mm,top=2.3mm,bottom=2.3mm,fonttitle=\bfseries,coltitle=white}}
% --- boîtes TUTORIEL (noms EXACTS, ne pas renommer) ---
\newtcolorbox{cle}[1][]{boxbase,colback=defcol!6,colframe=defcol,
  title={\textsc{À retenir}\ifx\relax#1\relax\relax\else\textmd{ : \itshape #1}\fi}}
\newtcolorbox{methode}[1][]{boxbase,colback=methcol!7,colframe=methcol,sharp corners,
  title={\textsc{Méthode}\ifx\relax#1\relax\relax\else\textmd{ --- \itshape #1}\fi}}
\newtcolorbox{exemplebox}[1][]{boxbase,colback=excol!5,colframe=excol,
  title={\textsc{Exemple}\ifx\relax#1\relax\relax\else\textmd{ : \itshape #1}\fi}}
\newtcolorbox{piege}[1][]{boxbase,colback=attncol!6,colframe=attncol,
  title={\textsc{Piège}\ifx\relax#1\relax\relax\else\textmd{ : \itshape #1}\fi}}
\newtcolorbox{remarque}[1][]{boxbase,colback=black!4,colframe=black!45,
  title={\textsc{Remarque}\ifx\relax#1\relax\relax\else\textmd{ : \itshape #1}\fi}}
% --- boîtes EXERCICE / CORRIGÉ (noms EXACTS, auto-comptées) ---
\newtcolorbox[auto counter]{exo}[1][]{boxbase,colback=primary!4,colframe=primary,
  title={\textsc{Exercice}~\thetcbcounter\ifx\relax#1\relax\relax\else\textmd{ --- \itshape #1}\fi}}
\newtcolorbox[auto counter]{corrige}[1][]{boxbase,colback=excol!4,colframe=excol,
  title={\textsc{Corrigé}~\thetcbcounter\ifx\relax#1\relax\relax\else\textmd{ --- \itshape #1}\fi}}
\newcommand{\vv}[1]{\vec{#1}}\newcommand{\ds}{\displaystyle}
\newcommand{\R}{\mathbb{R}}\newcommand{\N}{\mathbb{N}}\newcommand{\Z}{\mathbb{Z}}
% (un \usepackage{fancyhdr}+en-tête est possible ; il est ignoré côté web)
% =========================================================================
\begin{document}
\begin{exo}[champ d'un fil rectiligne]
Un long fil rectiligne est parcouru par un courant $I=\qty{10}{\ampere}$. On calcule le champ
magnétique à la distance $d=\qty{5}{\centi\metre}$ du fil.
\begin{center}
\begin{tikzpicture}[>=Stealth,scale=1.0]
  \draw[red!80!black,line width=2pt] (0,-1.3)--(0,1.3);
  \draw[red!80!black,->,line width=2pt] (0,0.3)--(0,0.9); \node[red!80!black,left] at (-0.05,1.0){$I$};
  \fill[black] (1.8,0) circle (1.3pt) node[below,font=\footnotesize]{$M$};
  \draw[<->,gray,line width=0.5pt] (0,-0.8)--(1.8,-0.8); \node[gray,below,font=\footnotesize] at (0.9,-0.82){$d$};
\end{tikzpicture}
\end{center}
Calcule l'intensité $B$ du champ au point $M$ (on prend $\mu_0/2\pi=2\times10^{-7}$).
\end{exo}

\begin{exo}[champ au centre d'une spire]
Une spire circulaire unique de rayon $r=\qty{2}{\centi\metre}$ est parcourue par un courant
$I=\qty{5}{\ampere}$. Calcule l'intensité du champ magnétique \textbf{au centre} de la spire
(on prend $\mu_0=4\pi\times10^{-7}\approx\qty{1.26e-6}{}$, et $2\pi\approx6{,}28$).
\end{exo}

\begin{exo}[champ dans un solénoïde]
Un solénoïde comporte $N=1000$ spires réparties sur une longueur $L=\qty{50}{\centi\metre}$ ;
il est parcouru par $I=\qty{3}{\ampere}$. Calcule l'intensité du champ magnétique à l'intérieur.
\end{exo}

\begin{exo}[reconnaître un électro-aimant]
On réalise le montage ci-dessous : une pile alimente, via un interrupteur, une bobine enroulée
autour d'un \textbf{clou en fer doux}.
\begin{center}
{\shorthandoff{;:!?}%
\begin{circuitikz}[european,scale=0.9]
  \draw (0,0) to[battery1, l=pile] (0,2)
        to[nos, l=inter.] (3.4,2)
        to[L] (3.4,0) -- (0,0);
  \node[right,font=\footnotesize,align=left] at (3.7,1){bobine autour\\ d'un clou en fer};
\end{circuitikz}}
\end{center}
\begin{enumerate}[label=\alph*)]
  \item Comment s'appelle ce dispositif lorsqu'un courant circule ?
  \item Que devient le champ magnétique si l'on \textbf{ouvre} l'interrupteur ? Justifie avec la
        nature du fer doux.
\end{enumerate}
\end{exo}

\begin{exo}[la décroissance en $1/d$]
À la distance $d=\qty{10}{\centi\metre}$ d'un fil rectiligne, le champ vaut $B=\qty{6e-5}{\tesla}$.
Sans recalculer $I$, en utilisant seulement la dépendance $B\propto 1/d$ :
\begin{enumerate}[label=\alph*)]
  \item Quelle est l'intensité du champ à $d'=\qty{20}{\centi\metre}$ ?
  \item Et à $d''=\qty{5}{\centi\metre}$ ?
\end{enumerate}
\end{exo}

\end{document}
